\begin{abstract}
AI 系统已经能够提出想法、修改代码、运行实验和撰写研究报告。更难的问题是：它们能否做出更好的科研决策，选择有用的未知问题，设计有信息量的检验，拒绝证据不足的解释，并把经验带到新的项目中？本文按照科研品味、问题形成、可靠执行、经验学习与评价这五个问题，综合 Awesome Autoresearch 集合中的 79 项资源。我们比较不同方法使用的反馈与证据，而不把完成科研流程当作科学能力的单一指标。集合包含 72 篇论文、6 篇文章或演讲以及 1 个软件项目。每个问题章节以局限、开放问题和两项具体实验方向结束，随后给出三个可优先启动的项目及共享评价协议。本文以 AI research 为主要 testbed，数学为补充场景。提出的实验均为研究设计，不是已经获得的结果，也不构成对方法新颖性的确认。
\end{abstract}

\noindent\textbf{关键词：}autoresearch；AI for meta-science；科研品味；假设生成；AI scientists；科研 agent 学习。
